\documentclass[book.tex]{subfiles}
%\newcommand{\ec}{\bf{r}_1, \ldots \bf{r}_N}
%\newcommand{\nc}{\bf{R}_1, \ldots \bf{R}_K}


\begin{document}
\chapter{Integral Programming}
\label{ch:integralprog}
\noindent\rule{11cm}{0.4pt} \\
\textbf{Prerequisites:}  \autoref{chap:ml_basics} %\autoref{chap:numerical_integration}, \autoref{chap:diffprog}  
\\
\textbf{Difficulty Level:} ** \\
\noindent\rule{11cm}{0.4pt}
\vspace{5mm}


Integral programs are differentiable programs which include integration as a construct. Many physical equations naturally feature integrals in their descriptions, making the ability to manipulate integrals within a differentiable program a powerful advantage. Let's start with a simple mathematical example. Assume that we want to learn a function $f$ with the form

\begin{align} f(x) = \int_0^x g(y, \theta) dy \end{align}

Here $\theta$ is a learnable parameter, $f$ is the function we'd like to learn and $g$ is its derivative. This general pattern arises often in physics, when we seek the solution to a set of differential equations. How can we optimize this function? Suppose that we have a dataset with training data points $X$ and labels $y$

\begin{align}
    X &= x_1,\dotsc, x_n \\
    y &= y_1, \dotsc, y_n
\end{align}

Assume that we're performing a regression task with $L^2$ loss. The learning task then is

\begin{align}
    \arg\min_{\theta} \sum_{i=1}^N \left \| \int_0^{x_i} g(y, \theta) dy - y_i\right \|^2
\end{align}

Differentiating the loss $\mathcal{L}$ through the integral sign, we'll be left with a term that looks like 

\begin{align}
\nabla \mathcal{L} \approx \int_0^x \frac{\partial}{ \partial \theta} g(y, \theta) dy \end{align}

Integrating arbitrary functions can be extremely challenging and is not possible to do in general except numerically. So, we fall back to some type of numerical solver. Assume that we have an algorithm $\textrm{Integrate}$ that implements a numerical solver. That is

\begin{align} 
\int_0^x g(y, \theta) dy \approx \textrm{Integrate}(g, x, \theta) 
\end{align}

Then we can rewrite our learning task as

\begin{align}
    \arg\min_{\theta} \sum_{i=1}^N \| \textrm{Integrate}(g, x_i, \theta) - y_i\|^2
\end{align}

If $\textrm{Integrate}$ is a standard numerical integration scheme, it is easy to show that $\textrm{Integrate}$ is a differentiable function. This means that every term in the loss equation above is differentiable. We have succeeded in reducing an integral program into a differentiable program.

In practice, the choice of integration scheme for $\textrm{Integrate}$ will be critical. The example we covered above is 1-dimensional, but in practice, the integral program we seek to solve might involve a set of high dimensional differential equations. More sophisticated Monte Carlo integration schemes are required to solve such integrals.

\section{Physical Examples}
We consider a few examples of integral programs in modeling physical systems.

\subsection{Variational Problems}

Integral programs naturally lend themselves to solving variational problems which arise in physics. For example, consider the Euler-Lagrange equation

\begin{align}
    S[q] &= \int_{x_1}^{x_2} L(x, q(x), q'(x)) dx
\end{align}

Here $S$ is a quantity commonly known as the action functional and $L$ is the Lagrangian. The minima of the action functional provides the equations of motion for a physical system. We can formulate the action function as an integral program

\begin{lstlisting}[language=python]
def S(q: [$x_1$, $x_2$] $\to$ $\mathbb{R}$) $\to$ $\mathbb{R}$:
  return Integrate(q, $x_1$, $x_2$, L(x, q, $\frac{d}{dt}$(q))
\end{lstlisting}

This differentiable program can be optimized directly by choosing a suitable parameterized form of \lstinline{q} (perhaps as a fully connected network).

\subsection{Path Integral Programming}

One of the most powerful tools of mathematical physics which we will develop in this book is the path integral. Path integral calculations are fundamental to modern quantum mechanics and modern quantum field theory. In this section, we will introduce some of their mathematical foundations and explain how we can use path integrals in integral programs. The foundation for path integral calculations is evaluation of Gaussian integrals so we will start by considering some Gaussian integrals. These integrals will be of considerable interest when we return to quantum field theory later in this book.

\subsection{Gaussian Integrals}


Let's now consider a very simple integral program

\begin{lstlisting}[language=python]
def gaussian_integral(a: $\mathbb{R}$, J: $\mathbb{R}$):
  return $\int_{-\infty}^{\infty}$ exp(-$\frac{1}{2}$a*x$^2$ + J*x)
\end{lstlisting}

We can have Physika numerically evaluate this integral. This integral also has a direct value which can be computed by completing the square so the numerical evaluation isn't useful in practice.

We can perturb this integral to consider 
\begin{align} G = \int_{-\infty}^{\infty} e^{-\frac{1}{2}ax^2 + iJ x} dx \end{align}

where $i$ is the imaginary number. Our strategy of completing the square works here so we can simply plug into a standard formula and obtain.

\begin{align} G = \sqrt{\frac{2 \pi}{a}} e^{(iJ)^2/2a} = \sqrt{\frac{2 \pi}{a}} e^{-J^2/2a} \end{align}

Here's one more variant with $a$ replaced by $-ia$

\begin{align} \int_{-\infty}^{\infty} e^{\frac{1}{2}iax^2 + iJx} dx \end{align}

Let's again plug into the formula we derived above.

\begin{align} G = \sqrt{\frac{2 \pi}{-ia}} e^{(iJ)^2/(-2ia)} = \sqrt{\frac{2 \pi i}{a}} e^{-iJ^2/2a} \end{align}

What happens though when the integral we consider becomes non-Gaussian?

\begin{align} H = \int_{-\infty}^{\infty} e^{-\frac{1}{2}ax^2 + b x^4} dx \end{align}
An integral program can be defined as

\begin{lstlisting}[language=python]
def nongaussian_integral(a: $\mathbb{R}$, J: $\mathbb{R}$):
  return $\int_{-\infty}^{\infty}$ exp(-$\frac{1}{2}$a*x$^2$ + J*x$^4$)
\end{lstlisting}
and evaluated numerically. Techniques from perturbation theory can be used to perform approximations in a structured fashion which we can encode into algorithmic solvers.
\end{document}